\documentclass[12pt,a4paper]{article}
\usepackage{geometry}
\geometry{margin=2.54cm}
\usepackage{graphicx}
\usepackage{booktabs}
\usepackage{grffile}
\usepackage{xcolor}
\usepackage{hyperref}
\hypersetup{colorlinks=true, linkcolor=blue, urlcolor=blue}
\usepackage{xeCJK}
\setCJKmainfont{SimSun}
\setCJKsansfont{SimHei}
\setCJKmonofont{FangSong}
\title{000002 万科A 多智能体研究报告}
\author{FinAgent}
\date{\today}
\begin{document}
\maketitle
\tableofcontents
\newpage
\sloppy
\section{000002 万科A 多智能体研究报告}

> **草稿状态**：验证评分 65，结论 `revise`。 本报告尚未通过验证 Agent 复核，请勿对外发布。

\subsection{ 目录}

\begin{itemize}
\item [执行摘要](\#执行摘要)
\item [核心指标速览](\#核心指标速览)
\item [量价与技术面](\#量价与技术面)
\item [基本面与估值](\#基本面与估值)
\item [资金与交易结构](\#资金与交易结构)
\item [宏观利率背景](\#宏观利率背景)
\item [综合风险与数据局限](\#综合风险与数据局限)
\item [数据与工具说明](\#数据与工具说明)
\item [免责声明](\#免责声明)
\end{itemize}

<a id=''执行摘要''></a>

\subsection{ 执行摘要}
万科A最新收盘价3.55元，近20日下跌8.27\%、近60日跌幅达23.98\%，技术面RSI仅26.72已进入超卖区，但基本面持续承压——2025年归母净利润亏损88.56亿元，市净率仅0.28倍，资产负债率仍高达77.13\%。5月29日单日换手率骤升至4.65\%、成交额16.13亿元，同时两融余额约25.78亿元，显示多空分歧加剧。

<a id=''核心指标速览''></a>

\subsection{ 核心指标速览}
\begin{table}[htbp]
\centering
\begin{tabular}{lr}
\toprule
指标 & 数值 \\
\midrule
中信一级行业 & 42 \\
最新收盘价 & 3.55 \\
MA20 & 3.7295 \\
MA60 & 4.027 \\
20 日收益率 & -8.27\% \\
60 日收益率 & -23.98\% \\
RSI14 & 26.7241 \\
20 日均量 & 1.50 亿 \\
PE(TTM) & -0.4799 \\
PB(TTM) & 0.2849 \\
PS(TTM) & 0.1888 \\
股息率(TTM) & 0.00\% \\
总市值 & 423.54 亿 \\
融资余额 & 25.72 亿 \\
融资买入额 & 3598.03 万 \\
\bottomrule
\end{tabular}
\end{table}


<a id=''量价与技术面''></a>

\subsection{ 量价与技术面}

**核心结论**

截至2026-05-29，万科A（000002.XSHE）收盘价3.55元，较60日均线折价11.85\%，RSI 14为26.72，处于超卖区域，但近20日跌幅达8.27\%，量价背离明显。

\#\#\#\# 图 · RSI 与 MACD

\begin{figure}[htbp]
\centering
\includegraphics[width=0.8\textwidth]{charts/000002_multi_agent_report/technical_indicators.png}
\caption{technical\_indicators.png}
\end{figure}

**图注** RSI 处于超卖区域；MACD 位于信号线下方，动能偏弱。

\#\#\#\# 表 · 技术指标快照

\begin{table}[htbp]
\centering
\begin{tabular}{llllllll}
\toprule
维度 & 最新收盘价 & MA20 & MA60 & 20 日收益率 & 60 日收益率 & RSI14 & 20 日均量 \\
\midrule
最新 & 3.55 & 3.7295 & 4.027 & -8.27\% & -23.98\% & 26.7241 & 1.50 亿 \\
\bottomrule
\end{tabular}
\end{table}


**趋势概览**

\#\#\#\# 图 · 收盘价与 MA20/MA60

\begin{figure}[htbp]
\centering
\includegraphics[width=0.8\textwidth]{charts/000002_multi_agent_report/moving_averages.png}
\caption{moving\_averages.png}
\end{figure}

**图注** 价格运行在均线之下，短中期趋势偏弱。

截至2026-05-29，万科A（000002.XSHE）收盘价3.55元，已跌破MA20（3.73元）和MA60（4.03元），呈现空头排列。近20日收益率为-8.27\%，近60日收益率深跌至-23.98\%，显示中期下行趋势持续。RSI 14为26.72，进入技术性超卖区间（通常低于30），表明短期抛压可能过度，但尚未出现明确反转信号。

\#\#\#\# 图 · 回撤曲线

\begin{figure}[htbp]
\centering
\includegraphics[width=0.8\textwidth]{charts/000002_multi_agent_report/drawdown.png}
\caption{drawdown.png}
\end{figure}

**图注** 回撤经历加深后有所修复。

\begin{table}[htbp]
\centering
\begin{tabular}{lll}
\toprule
指标 & 数值 & 日期 \\
\midrule
最新收盘价 & 3.55元 & 2026-05-29 \\
MA20 & 3.73元 & 2026-05-29 \\
MA60 & 4.03元 & 2026-05-29 \\
20日收益率 & -8.27\% & 2026-05-29 \\
60日收益率 & -23.98\% & 2026-05-29 \\
RSI 14 & 26.72 & 2026-05-29 \\
\bottomrule
\end{tabular}
\end{table}


**近期异动**

2026-05-29，万科A（000002.XSHE）股价从前一日低点3.31元大幅反弹至3.55元，单日涨幅7.25\%，成交量达4.51亿股，是20日均量（1.50亿股）的3.01倍，换手率显著放大。此前，股价于2026-05-28创下区间最低收盘价3.31元，当日成交量1.21亿股，低于均量，显示下跌缩量特征。

\begin{table}[htbp]
\centering
\begin{tabular}{lllll}
\toprule
日期 & 收盘价 & 涨跌幅 & 成交量 & 较20日均量倍数 \\
\midrule
2026-05-28 & 3.31 & -0.60\% & 1.21亿 & 0.81x \\
2026-05-29 & 3.55 & +7.25\% & 4.51亿 & 3.01x \\
\bottomrule
\end{tabular}
\end{table}


**量价关系**

\#\#\#\# 图 · 收盘价与成交量

\begin{figure}[htbp]
\centering
\includegraphics[width=0.8\textwidth]{charts/000002_multi_agent_report/price_volume.png}
\caption{price\_volume.png}
\end{figure}

**图注** 价格整体下行，成交量先抑后扬。

近20日（2026-04-29至2026-05-29），万科A（000002.XSHE）股价从3.87元跌至3.55元，跌幅8.27\%，而同期成交量从1.45亿股增至4.51亿股，增幅210.89\%，呈现典型的**放量下跌**形态，表明下跌过程中抛压与承接均较活跃。然而，在2026-05-29的反弹中，成交量急剧放大至均量的3倍，形成**放量反弹**，可能暗示短期资金博弈加剧。结合RSI超卖，量价关系显示市场在低位分歧加大，但趋势尚未逆转。

**融资余额变化**

万科A（000002.XSHE）融资余额从2025-09-11的34.63亿元下降至2026-05-28的25.72亿元，降幅25.73\%，期间峰值出现在2025-12-10的4.09亿元。融资余额持续收缩，反映杠杆资金在股价下行过程中逐步离场，与价格趋势方向一致。

**数据局限**

\begin{itemize}
\item 未提供融券余额数据，无法评估融券做空力量。
\item 未提供大宗交易、龙虎榜等资金流向数据，无法判断主力资金动向。
\item 未提供波动率指标（如ATR、布林带宽度），无法量化价格波动幅度。
\end{itemize}

<a id=''基本面与估值''></a>

\subsection{ 基本面与估值}

**核心结论**

截至2026-05-29，万科A连续三年亏损且经营现金流由正转负，PE为负、PB仅0.28倍，盈利质量与偿债能力持续恶化。

\#\#\#\# 表 · 最新偿债与流动性

\begin{table}[htbp]
\centering
\begin{tabular}{llll}
\toprule
维度 & 流动比率 & 速动比率 & 资产负债率 \\
\midrule
最新 & 1.22 倍 & 0.56 倍 & 77.13\% \\
\bottomrule
\end{tabular}
\end{table}


**盈利与成长**

\#\#\#\# 表 · 最新盈利质量因子

\begin{table}[htbp]
\centering
\begin{tabular}{llll}
\toprule
维度 & 毛利率(TTM) & 净利率(TTM) & ROE(TTM) \\
\midrule
最新 & 9.99\% & -40.78\% & — \\
\bottomrule
\end{tabular}
\end{table}


2023至2025年，万科A（000002.XSHE）营收与利润持续下滑，归母净利润由盈转亏且亏损逐年扩大。毛利率从2023年的11.1\%降至2025年的9.5\%，净利率从2.6\%骤降至-40.8\%。MD\&A指出，亏损主因房地产开发项目结算规模下降、毛利率低位、信用及资产减值计提增加。如图1所示，营收增速与毛利率同步下行，反映行业下行周期中公司盈利能力的系统性恶化。

\begin{table}[htbp]
\centering
\begin{tabular}{llll}
\toprule
指标 & 2023年 & 2024年 & 2025年 \\
\midrule
营收（亿元） & 4657.39 & 3431.76 & 2334.33 \\
营收同比增速 & - & -26.3\% & -32.0\% \\
归母净利润（亿元） & 121.63 & -494.78 & -885.56 \\
毛利率（TTM） & 11.1\% & 10.2\% & 9.5\% \\
净利率（TTM） & 2.6\% & -14.4\% & -40.8\% \\
\bottomrule
\end{tabular}
\end{table}


**现金流与勾稽**

2025年经营现金流由正转负，收现比仅0.53，净现比接近0，利润与现金流严重背离。MD\&A解释，经营活动现金净流出9.9亿元，主因销售回款减少及流动性压力加大。应收账款增速（16.9\%）快于收入增速（-32.0\%），回款节奏滞后。如图2所示，经营现金流与资本开支的缺口扩大，进一步加剧流动性压力。

\begin{table}[htbp]
\centering
\begin{tabular}{llll}
\toprule
指标 & 2023年 & 2024年 & 2025年 \\
\midrule
经营现金流（亿元） & 39.12 & 38.00 & -9.88 \\
收现比 & 0.84 & 0.64 & 0.53 \\
净现比 & 0.32 & -0.08 & 0.01 \\
\bottomrule
\end{tabular}
\end{table}


**估值水平**

\#\#\#\# 图 · 估值因子走势

\begin{figure}[htbp]
\centering
\includegraphics[width=0.8\textwidth]{charts/000002_multi_agent_report/valuation_factors.png}
\caption{valuation\_factors.png}
\end{figure}

**图注** 多条曲线整体向下，指标同步走弱。

\#\#\#\# 图 · PE/PB 历史分位

\begin{figure}[htbp]
\centering
\includegraphics[width=0.8\textwidth]{charts/000002_multi_agent_report/valuation_percentile.png}
\caption{valuation\_percentile.png}
\end{figure}

截至2026-05-29，PE（TTM）为-0.48倍（因亏损），PB为0.28倍，PS为0.19倍，股息率为0。对比2025-09-11，PB从0.40倍降至0.28倍，市值从814.87亿元缩水至423.54亿元，反映市场对资产价值持续折价。如图3所示，PE、PB、PS均处于历史低位，且低于历史百分位10\%以下，显示市场对房地产行业及公司基本面的悲观预期。

\begin{table}[htbp]
\centering
\begin{tabular}{lll}
\toprule
估值指标 & 2025-09-11 & 2026-05-29 \\
\midrule
PE（TTM） & -1.58 & -0.48 \\
PB & 0.40 & 0.28 \\
PS & 0.27 & 0.19 \\
股息率 & 0.0\% & 0.0\% \\
市值（亿元） & 814.87 & 423.54 \\
\bottomrule
\end{tabular}
\end{table}


\begin{itemize}
\item 数据局限：未提供2023年及之前的完整财务数据，无法计算更长周期的增长率；无券商预测数据，无法评估市场一致预期。
\end{itemize}

<a id=''资金与交易结构''></a>

\subsection{ 资金与交易结构}

**核心结论**

截至2026年5月29日，万科A融资余额较2025年9月11日下降25.7\%至25.72亿元，同期股价下跌49.6\%，资金持续流出且杠杆去化显著。

\#\#\#\# 图 · 融资融券余额

\begin{figure}[htbp]
\centering
\includegraphics[width=0.8\textwidth]{charts/000002_multi_agent_report/margin_balances.png}
\caption{margin\_balances.png}
\end{figure}

**图注** 融资余额曲线整体下行。

\#\#\#\# 表 · 两融区间变动

\begin{table}[htbp]
\centering
\begin{tabular}{ll}
\toprule
指标 & 数值 \\
\midrule
区间起始 & 2025-09-11 \\
区间结束 & 2026-05-28 \\
融资余额（期初） & 34.63 亿 \\
融资余额（期末） & 25.72 亿 \\
融资余额变动 & -8.91 亿 \\
变动幅度 & -25.73\% \\
融券余额（期初） & 0.19 亿 \\
融券余额（期末） & 0.06 亿 \\
融券余额变动 & -0.13 亿 \\
融券变动幅度 & -66.47\% \\
融资买入峰值 & 4.09 亿（2025-12-10） \\
\bottomrule
\end{tabular}
\end{table}


**成交与活跃度**

\#\#\#\# 表 · 成交活跃度

\begin{table}[htbp]
\centering
\begin{tabular}{ll}
\toprule
指标 & 数值 \\
\midrule
20 日均量 & 1.50 亿 \\
近12日日均成交额 & 5.29 亿 \\
近12日最高成交额 & 16.13 亿（2026-05-29） \\
近12日最低成交额 & 3.01 亿（2026-05-22） \\
近12日换手率区间 & 89.18\% \textasciitilde{} 4.65\% \\
主力资金流向 & 数据缺失 \\
\bottomrule
\end{tabular}
\end{table}


2026年5月25日至29日（最近5个交易日），万科A（000002.XSHE）成交额从3.30亿元激增至16.13亿元，增幅达388.5\%；同期收盘价从3.49元微涨至3.55元，涨幅1.7\%。更长时间维度看，最近20个交易日（2026年4月29日至5月29日）成交额从4.68亿元升至16.13亿元，增幅244.9\%，而股价下跌8.3\%，呈现放量下跌态势（如图1所示）。区间内（2025年9月11日至2026年5月29日），股价最高点出现在2025年9月17日（7.08元），最低点出现在2026年5月28日（3.27元），区间振幅达116.5\%。

截至2026年5月29日，RSI14为26.7，处于超卖区域。

\begin{table}[htbp]
\centering
\begin{tabular}{lllll}
\toprule
窗口 & 起始日 & 结束日 & 收盘价变化 & 成交额变化 \\
\midrule
最近5日 & 2026-05-25 & 2026-05-29 & +1.7\% & +388.5\% \\
最近20日 & 2026-04-29 & 2026-05-29 & -8.3\% & +244.9\% \\
\bottomrule
\end{tabular}
\end{table}


**融资融券**

\#\#\#\# 图 · 融资余额与买卖额

\begin{figure}[htbp]
\centering
\includegraphics[width=0.8\textwidth]{charts/000002_multi_agent_report/margin_enhanced.png}
\caption{margin\_enhanced.png}
\end{figure}

\#\#\#\# 表 · 融资融券快照

\begin{table}[htbp]
\centering
\begin{tabular}{lllllll}
\toprule
统计日期 & 融资余额 & 融券余额 & 两融余额 & 融资买入额 & 融资偿还额 & 融券余量 \\
\midrule
2026-05-28 & 25.72 亿 & 0.06 亿 & 25.78 亿 & 0.36 亿 & 0.54 亿 & 191.02 万 \\
\bottomrule
\end{tabular}
\end{table}


融资余额从2025年9月11日的34.63亿元持续下降至2026年5月28日的25.72亿元，区间降幅25.7\%（如图2所示）。融资买入峰值出现在2025年12月10日，当日融资买入额为4.09亿元。融券余额数据缺失（JSON中`short\_selling\_balance\_end`为null），无法评估融券端变化。

\begin{table}[htbp]
\centering
\begin{tabular}{llllll}
\toprule
指标 & 起始日 & 数值 & 结束日 & 数值 & 变化 \\
\midrule
融资余额 & 2025-09-11 & 34.63亿元 & 2026-05-28 & 25.72亿元 & -25.7\% \\
峰值融资买入 & - & - & 2025-12-10 & 4.09亿元 & - \\
\bottomrule
\end{tabular}
\end{table}


**股东与股本**

\#\#\#\# 表 · 股本结构快照

\begin{table}[htbp]
\centering
\begin{tabular}{ll}
\toprule
项目 & 数值 \\
\midrule
统计日期 & 2026-05-29 \\
总股本 & 119.31 亿 \\
流通 A 股 & 97.17 亿 \\
自由流通股本 & 64.74 亿 \\
自由流通占比 & 54.3\% \\
非自由流通占比 & 45.7\% \\
\bottomrule
\end{tabular}
\end{table}


截至2026年5月29日，万科A（000002.XSHE）总市值423.54亿元，较2025年9月11日的814.87亿元下降48.0\%。市盈率（TTM）为-0.48倍（持续亏损），市净率（TTM）为0.28倍，市销率（TTM）为0.19倍。资产负债率从2025年9月11日的73.1\%升至2026年5月29日的77.1\%，杠杆水平持续攀升。

\begin{table}[htbp]
\centering
\begin{tabular}{llll}
\toprule
指标 & 2025-09-11 & 2026-05-29 & 变化 \\
\midrule
总市值 & 814.87亿元 & 423.54亿元 & -48.0\% \\
市净率 & 0.40倍 & 0.28倍 & -30.0\% \\
资产负债率 & 73.1\% & 77.1\% & +4.0个百分点 \\
\bottomrule
\end{tabular}
\end{table}


**分红与资金成本**

截至2026年5月29日，万科A（000002.XSHE）股息率（TTM）为0.0\%，表明最近12个月未进行现金分红。存量融资综合成本为3.02\%（MD\&A披露），较上年末下降85个基点。2025年实际利息支出合计129.9亿元，其中资本化利息51.6亿元。

**现金流与利润质量**

根据财务报表，2023-2025年营收、归母净利润与经营现金流呈现持续恶化趋势：

\begin{table}[htbp]
\centering
\begin{tabular}{llll}
\toprule
年份 & 营收（亿元） & 归母净利润（亿元） & 经营现金流（亿元） \\
\midrule
2023 & 4,657.39 & 121.63 & 39.12 \\
2024 & 3,431.76 & -494.78 & 38.00 \\
2025 & 2,334.33 & -885.56 & -9.88 \\
\bottomrule
\end{tabular}
\end{table}


MD\&A解释2025年经营现金流为净支出9.9亿元，主要因经营活动现金流入同比下降36.18\%至1,401.73亿元，而现金流出仅下降34.59\%至1,411.61亿元。收现比（经营现金流入/营收）为0.53，低于1，表明收入中有相当比例尚未转化为销售回款。利润增长与经营现金流变动背离（归母净利润同比改善79.0\%，但经营现金流同比恶化-126.0\%），利润质量存在风险。

**数据局限**
\begin{itemize}
\item 融券余额数据缺失，无法分析融券端交易行为。
\item 股东户数、前十大股东持股比例等股本结构数据未提供。
\item 分红历史数据仅提供TTM股息率，无法追溯过往分红政策变化。
\item 自由现金流数据仅提供2025年单年数值，无法进行趋势对比。
\end{itemize}

<a id=''宏观利率背景''></a>

\subsection{ 宏观利率背景}

**核心结论**

截至2026-05-29，万科A（000002.XSHE）股价3.55元，近20日下跌8.27\%，因JSON数据中无短端利率及国债收益率数据，无法评估利率环境对估值的传导。

**趋势概览**

**短端利率**
\begin{itemize}
\item 数据局限：JSON中未提供Shibor或任何短端利率数据，无法计算最新值及与20日前的变动。
\end{itemize}

\#\#\#\# 图 · Shibor 利率

\begin{figure}[htbp]
\centering
\includegraphics[width=0.8\textwidth]{charts/000002_multi_agent_report/shibor_rates.png}
\caption{shibor\_rates.png}
\end{figure}

**图注** 短端利率整体下行，波动相对平稳。

**收益率曲线**
\begin{itemize}
\item 数据局限：JSON中未提供1Y/10Y/30Y国债收益率或期限利差数据，无法进行曲线形态分析。
\end{itemize}

\#\#\#\# 图 · 国债收益率走势

\begin{figure}[htbp]
\centering
\includegraphics[width=0.8\textwidth]{charts/000002_multi_agent_report/gov_yield_trend.png}
\caption{gov\_yield\_trend.png}
\end{figure}

**图注** 多条曲线整体向下，指标同步走弱。

**利率与估值的逻辑联系**
\begin{itemize}
\item 数据局限：JSON中未提供无风险利率（如10Y国债收益率）数据，无法计算股权风险溢价或与股息率（当前为0.0\%）进行对比。因此，无法评估利率环境对万科A（000002.XSHE）估值的直接或间接影响。
\end{itemize}

**数据局限**

\begin{itemize}
\item 或股息率（0.0\%）之间的逻辑联系
\end{itemize}

<a id=''综合风险与数据局限''></a>

\subsection{ 综合风险与数据局限}

**核心结论**

截至2026-05-29，万科A股价3.55元，RSI 26.7进入超卖区，但归母净利润连续两年亏损且经营现金流由正转负，综合风险较高。

**价格与波动**
截至2026-05-29，万科A（000002.XSHE）收盘价3.55元，较60日均线（4.03元）折价11.8\%，RSI 14为26.7，已进入超卖区域（<30）。近60日累计下跌24.0\%，区间内股价从2025-09-17高点7.08元跌至2026-05-28低点3.27元，最大回撤达53.8\%。近5个交易日（2026-05-25至2026-05-29）股价反弹1.7\%，但成交量放大3.8倍（从9501万股增至4.51亿股），显示多空分歧加剧。

**基本面与估值**
万科A（000002.XSHE）估值指标显示深度亏损状态：PE（TTM）为-0.48倍，PB为0.28倍，PS为0.19倍。归母净利润连续两年亏损且亏损幅度扩大，2024年为-494.8亿元，2025年进一步扩大至-885.6亿元。毛利率从2023年的11.1\%降至2025年的10.0\%，净利率从2023年的-16.7\%恶化至2025年的-40.8\%。营收从2023年的4657.4亿元降至2025年的2334.3亿元，降幅达49.9\%。

MD\&A解释为“房地产开发项目结算规模显著下降，毛利率仍处低位”及“新增计提了信用减值和资产减值”。

\begin{table}[htbp]
\centering
\begin{tabular}{llll}
\toprule
指标 & 2023年 & 2024年 & 2025年 \\
\midrule
营收（亿元） & 4657.4 & 3431.8 & 2334.3 \\
归母净利润（亿元） & 121.6 & -494.8 & -885.6 \\
毛利率（TTM） & 11.1\% & 10.2\% & 10.0\% \\
净利率（TTM） & -16.7\% & -40.8\% & -40.8\% \\
\bottomrule
\end{tabular}
\end{table}


**杠杆与流动性**
万科A（000002.XSHE）资产负债率从2023年的73.1\%升至2025年的77.1\%，负债合计7847.6亿元。流动比率为1.22，速动比率为0.56，短期偿债能力偏弱。经营现金流从2024年的38.0亿元转为2025年的-9.9亿元，同比变动-126.0\%。MD\&A确认“全年经营活动产生的现金流量为净支出9.9亿元”，并提及“因面临公开债务的集中兑付，本集团流动性压力进一步加大”。融资余额从2025-09-11的34.6亿元降至2026-05-28的25.7亿元，降幅25.7\%，显示杠杆资金持续离场。

**宏观与利率**
无直接宏观或利率数据可用。MD\&A提及“存量融资的综合成本为3.02\%，较上年末下降85个基点”，显示融资成本有所下降，但未提供具体利率曲线或宏观环境分析。

**数据局限**

\begin{itemize}
\item ，无法评估股东回报
\end{itemize}

<a id=''数据与工具说明''></a>

\subsection{ 数据与工具说明}
\begin{itemize}
\item 数据区间：2025-09-11 至 2026-05-29
\item 计划使用的米筐函数：get\_price, get\_price\_change\_rate, get\_turnover\_rate, get\_capital\_flow, get\_factor, get\_securities\_margin, get\_dividend, get\_shares, get\_instrument\_industry, is\_suspended, is\_st\_stock, get\_interbank\_offered\_rate, get\_yield\_curve, get\_pit\_financials\_ex
\item 数据质量：9 项可用；缺失 基准指数, 资金流向, 大宗交易, 分红
\item 生成时间：2026-05-31 17:55
\end{itemize}

<a id=''免责声明''></a>

\subsection{ 免责声明}
本报告仅供课程研究与信息展示，不构成投资建议。\end{document}
